%%%%%%%%%%%%%%%%%%%%%%%%%%%%%%%%%%%%%%%%%%%%%%%%%%%%%%%%%%%%%%%%%%%%%%
% Lớp 10 | Hệ thức lượng trong tam giác | Giá trị lượng giác của góc (0 −180∘)
% Dạng 0H4?1-2: Tính các giá trị lượng giác
% Số câu: 29  (sinh tự động bởi tach_cau.py; từ nay thầy sửa trực tiếp file này)
%%%%%%%%%%%%%%%%%%%%%%%%%%%%%%%%%%%%%%%%%%%%%%%%%%%%%%%%%%%%%%%%%%%%%%

%%%%%%%%%%%%%%%% TRẮC NGHIỆM 4 PHƯƠNG ÁN %%%%%%%%%%%%%%%%

% nguồn: MaHoaNH_TracNghiem_Lop10
\begin{ex}%[0H4N1-2]%[ma:70346a81]
	 Trong các khẳng định sau, khẳng định nào đúng?
	 \choice
		{\True $\sin  30^\circ=\cos 60^\circ$}
		{$\sin 120^\circ=\cos 60^\circ$} 
		{$\tan 50^\circ=-\cot  40^\circ$}
		{$\sin  70^\circ=\sin 20^\circ$}
	 \loigiai{
	 Ta có: $\sin 30^\circ=\cos 60^\circ$.}
\end{ex}

% nguồn: MaHoaNH_TracNghiem_Lop10
\begin{ex}%[0H4N1-2]%[ma:c2a6ba8d]
	Giá trị của biểu thức  $A=\sin^2 51^\circ+\sin^2 55^\circ+\sin^2 39^\circ+\sin^2 35^\circ$ là
	 \choice
		 {\True $2$}
		 {$3$}
		 {$1$}
		 {$4$}
	\loigiai{ 
	$A=\left(\sin^2 51^\circ+\sin^2 39^\circ\right)+\left(\sin^2 55^\circ+\sin^2 35^\circ\right)=\left(\sin^2 51^\circ+\cos^2 51^\circ\right)+\left(\sin^2  55^\circ+\cos^2 55^\circ\right)=2$.}
\end{ex}

% nguồn: MaHoaNH_TracNghiem_Lop10
\begin{ex}%[0H4N1-2]%[ma:e19471e5]
	Giá trị của biểu thức $A=\tan 1^\circ\cdot\tan 2^\circ\cdot\tan 3^\circ\cdots\tan 88^\circ\cdot\tan  89^\circ$ là
	 \choice
		{$0$} 
		{$2$} 
		 {$3$}
		 {\True $1$}
	\loigiai{ 
	Ta có  $A=\left(\tan 1^\circ\cdot\tan 89^\circ\right)\cdot\left(\tan 2^\circ\cdot\tan 88^\circ\right)\cdots\left(\tan 44^\circ\cdot\tan 46^\circ\right)\cdot\tan 45^\circ=1$.}
\end{ex}

% nguồn: MaHoaNH_TracNghiem_Lop10
\begin{ex}%[0H4N1-2]%[ma:0f0042ad]
	Trong các khẳng định sau, khẳng định  nào \textbf{sai}?
	 \choice
		 {$\sin 0^\circ+\cos 0^\circ=1$}
		{$\sin  180^\circ+\cos 180^\circ=-1$}
		{$\sin  90^\circ+\cos 90^\circ=1$}
		{\True $\sin 60^\circ+\cos  60^\circ=1$}
	\loigiai{ 
	Ta  có $\sin 60^\circ+\cos 60^\circ=\dfrac{\sqrt{3}+1}{2}$.}
\end{ex}

% nguồn: MaHoaNH_TracNghiem_Lop10
\begin{ex}%[0H4N1-2]%[ma:710c42fd]
	Trong các đẳng  thức sau, đẳng thức nào đúng?
	 \choice
		{$\sin \left(180^\circ-\alpha  \right)=-\cos \alpha$}
		{$\sin \left(180^\circ-\alpha \right)=-\sin  \alpha$}
		 {\True $\sin \left(180^\circ-\alpha \right)=\sin \alpha$}
		 {$\sin \left(180^\circ-\alpha \right)=\cos \alpha$}
	\loigiai{ 
		Ta có $\sin \left(180^\circ-\alpha  \right)=\sin \alpha$.
	 }
\end{ex}

% nguồn: MaHoaNH_TracNghiem_Lop10
\begin{ex}%[0H4N1-2]%[ma:6737e8b7]
	Giá trị của $\cos 60^\circ+\sin 30^\circ$ bằng bao nhiêu? 
	 \choice
		 {$\dfrac{\sqrt{3}}{2}$}
		{$\sqrt{3}$} 
		{$\dfrac{\sqrt{3}}{3}$} 
		{\True  $1$}
	\loigiai{ 
		Ta có $\cos  60^\circ+\sin 30^\circ=\dfrac{1}{2}+\dfrac{1}{2}=1$.
	} 
\end{ex}

% nguồn: MaHoaNH_TracNghiem_Lop10
\begin{ex}%[0H4H1-2]%[ma:d8c19b46]
	Tính giá trị của biểu thức $A=\sin (10^{\circ})\cdot \sin (20^{\circ})\ldots \sin (190^{\circ})\cdot \sin (200^{\circ})$. 
	\choice 
		{$4$} 
		{$2$} 
		 {\True $0$}
		{$-2$} 
	 \loigiai{
		Ta thấy trong tích có $\sin 180^{\circ}=0$ và do đó  $A=0$.
	} 
\end{ex}

% nguồn: MaHoaNH_TracNghiem_Lop10
\begin{ex}%[0H4Y1-2]%[ma:97cfa209]
	Đường thẳng $d$ đi  qua điểm $A\left(-2;-3\right)$ và có VTCP $\overrightarrow{u}=\left(-2;1\right)$ có phương trình là
	\choice 
		{$\heva{&x=-2-2t\\&y=1-3t}$} 
		{$\heva{&x=-2-3t\\&y=1-2t}$} 
		{$\heva{&x=-2+t\\&y=-3-2t}$} 
		{\True $\heva{&x=-2-2t\\&y=-3+t}$} 
	\loigiai{Đường thẳng $d$ đi qua điểm $A\left(-2;-3\right)$ và có VTCP $\overrightarrow{u}=\left(-2;1\right)$ nên có phương trình  là$\colon\heva{&x=-2-2t\\&y=-3+t}$
	} 
\end{ex}

% nguồn: MaHoaNH_TracNghiem_Lop10
\begin{ex}%[0H4B1-2]%[ma:ec70fe80]
	Trong mặt phẳng $Oxy$, phương  trình tham số của đường thẳng $\Delta$ đi qua điểm $A(-4; 2)$ nhận $\vec{u}=(-2; 3)$ làm véc-tơ chỉ phương là
	 \choice
		 {\True $\heva{&x=-4-2t\\ &y=2+3t}$}
		 {$\heva{&x=4+2t\\ &y=-2-3t}$}
		{$\heva{&x=-2-4t\\  &y=3+2t}$}
		{$\heva{&x=-2-2t\\  &y=-1+3t}$}
	\loigiai{Phương trình tham số của đường thẳng $\Delta$ đi qua điểm $A(-4; 2)$ nhận  $\vec{u}=(-2; 3)$ làm véc-tơ chỉ phương là $\heva{&x=-4-2t\\ &y=2+3t.}$
	 }
\end{ex}

% nguồn: MaHoaNH_TracNghiem_Lop10
\begin{ex}%[0H4Y1-2]%[ma:2e4b9144]
	Trong mặt phẳng $Oxy$, phương trình tham số của đường thẳng $d$  đi qua điểm $M(1;-2)$ và có véc-tơ chỉ phương $\vec{u}=(3;5)$ là
	\choice 
		{$d\colon\heva{&x=1+5t\\&y=-2-3t}$} 
		{$d\colon\heva{&x=3+t\\&y=5-2t}$} 
		 {$d\colon\heva{&x=3+2t\\&y=5+t}$}
		{\True $d\colon\heva{&x=1+3t\\&y=-2+5t}$} 
	 \loigiai{
		Đường thẳng $d$ đi qua điểm $M(1;-2)$ và  có véc-tơ chỉ phương $\vec{u}=(3;5)$ có phương trình tham số $d\colon\heva{&x=1+3t\\&y=-2+5t.}$
	} 
\end{ex}

% nguồn: MaHoaNH_TracNghiem_Lop10
\begin{ex}%[0H4B1-2]%[ma:c8670623]
	Phương trình tham số của đường  thẳng đi qua hai điểm $A(2;-1)$ và $B(2;5)$ là
	\choice 
		{$\heva{&x=2t  \\ &y=-6t}$}
		{$\heva{&x=2+t \\ &y=  5+6t}$}
		{$\heva{&x=1 \\ &y=2+6t}$} 
		{\True  $\heva{&x=2 \\&y=-1+6t}$}
	 \loigiai{
		 Đường thẳng $AB$ có một véc tơ chỉ phương là $\overrightarrow{AB}=(0;6)$ và đi qua điểm $A(2;-1)$ nên có phương trình là $ \heva{&x=2 \\&y=-1+6t.}$
	} 
\end{ex}

% nguồn: MaHoaNH_TracNghiem_Lop10
\begin{ex}%[0H4Y1-2]%[ma:55cb4e8a]
	Trong mặt phẳng  $O x y$, tìm phương trình đường thẳng đi qua điểm $M(2; 0)$ và có một véc-tơ pháp tuyến $\vec{n}=(1;-2)$.
	 \choice
		{$x+2 y-2=0$} 
		{\True $x-2 y-2=0$} 
		{$x-2 y+2=0$} 
		{$2  x+y-4=0$}
	\loigiai{Phương trình đường thẳng đi qua điểm $M(2; 0)$ và có một véc-tơ pháp tuyến $\overrightarrow{n}=(1;-2)$  là
		 $$1\cdot(x-2)-2\cdot(y-0)=0\Leftrightarrow x-2y-2=0.$$
	} 
\end{ex}

% nguồn: MaHoaNH_TracNghiem_Lop10
\begin{ex}%[0H4B1-2]%[ma:eb445a3e]
	Cho $\triangle ABC$ có $A(2;-1)$; $B(4;5)$; $C(-3;2)$.  Viết phương trình tổng quát của đường cao $AH$.
	\choice 
		{\True $7x+3y-11=0  $}
		 {$3x+7y+1=0$}
		{$7x+3y+11=0$} 
		{$-7x+3y+11=0$} 
	\loigiai{ 
		$AH$ qua  $A(2;-1)$, có một véc-tơ pháp tuyến là $\overrightarrow{BC}=(-7;-3)$.\\
		Phương trình đường cao $AH$ là  $$ -7(x-2)-3(y+1)=0\Leftrightarrow -7x-3y+11=0\Leftrightarrow 7x+3y-11=0.$$
	} 
\end{ex}

% nguồn: MaHoaNH_TracNghiem_Lop10
\begin{ex}%[0H4B1-2]%[ma:cc5e9272]
	Trong mặt phẳng $Oxy$, phương trình  tổng quát của đường thẳng đi qua $A(2;-1)$, $B(3;-2)$ là
	 \choice
		{$x-y-3=0$} 
		{$x-y-1=0$} 
		 {$x+y+3=0$}
		{\True  $x+y-1=0$}
	\loigiai{ 
		Ta có $\overrightarrow{AB}=(1;-1)$, chọn $\overrightarrow{n}=(1;1)$ là một véc-tơ pháp tuyến  của đường thẳng $AB$.\\
		Phương trình tổng quát của đường thẳng $AB$  là $x+y-1=0$.
	} 
\end{ex}

% nguồn: MaHoaNH_TracNghiem_Lop10
\begin{ex}%[0H4Y1-2]%[ma:be2fc75a]
	Phương trình nào sau đây là phương trình tổng quát  của đường thẳng?
	 \choice
		{$y^2=4x$} 
		{$\heva{&x=4+3 t\\  &y=2-t}$}
		 {$\dfrac{x^2}{25}+\dfrac{y^2}{9}=1$}
		 {\True $4x-3y-2=0$}
	\loigiai{ 
		Phương trình tổng quát  của đường thẳng có dạng $ax+by+c=0$.\\
		Do đó đáp án đúng là phương trình  $4x-3y-2=0$.
	} 
\end{ex}

% nguồn: MaHoaNH_TracNghiem_Lop10
\begin{ex}%[0H4B1-2]%[ma:90e9e5e9]
	Trong mặt phẳng $Oxy$ cho điểm  $M(1;2)$. Gọi $A$, $B$ là hình chiếu của $M$ lên $Ox$, $Oy$. Viết phương trình đường thẳng $AB$.
	\choice 
		 {$x+y-3=0$}
		 {\True $2x+y-2=0$}
		{$x+2y-1=0$} 
		 {$2x+y+2=0$}
	 \loigiai{
	Ta có  $A(1;0)$ và $B(0;2)$. \\
	Phương trình đường thẳng $AB$  là $\dfrac{x}{1}+\dfrac{y}{2}=1 \Leftrightarrow 2x+y-2=0$.
	 }
\end{ex}

% nguồn: MaHoaNH_TracNghiem_Lop10
\begin{ex}%[0H4Y1-2]%[ma:f17ccfb3]
	Véc-tơ nào dưới  đây là $1$ véc-tơ chỉ phương của đường thẳng song song với trục $Ox$?
	\choice 
		 {$\overrightarrow{u}=(1;-1)$}
		 {$\overrightarrow{u}=(0;1)$}
		{$\overrightarrow{u}=(1;1)$} 
		{\True  $\overrightarrow{u}=(1;0)$}
	\loigiai{ 
	Trục $Ox$ có phương trình tổng quát là $y=0$, suy ra trục $Ox$  có $1$ véc-tơ chỉ phương là $(1;0)$. \\
	Vậy đường thẳng song song với trục $Ox$ có $1$ véc-tơ  chỉ phương là $\overrightarrow{u}=(1;0)$.
	 }
\end{ex}

% nguồn: MaHoaNH_TracNghiem_Lop10
\begin{ex}%[0H4B1-2]%[ma:668db43b]
	Phương trình tham số của  đường thẳng $d$ đi qua điểm $M(4;-7)$ và song song với trục hoành là
	\def\dotEX{} 
	\choice 
		 {$\heva{&x=4+t\\&y=-7.}$}
		{\True $\heva{&x=4\\&y=-7+t.}$} 
		{$\heva{&x=-7+t\\&y=4.}$} 
		 {$\heva{&x=1+4t\\&y=-7t.}$}
	\loigiai{ 
		Đường thẳng $d$ đi qua điểm $M(4;-7)$ và song song với trục hoành nhận $\vec{i}=(1;0)$ làm một véc-tơ chỉ phương, do đó $d$ có phương trình tham số  là $\heva{&x=4\\&y=-7+t.}$
	} 
\end{ex}

% nguồn: MaHoaNH_TracNghiem_Lop10
\begin{ex}%[0H4Y1-2]%[ma:56faa10f]
	Trong mặt phẳng $Oxy$, cho các điểm $M(2;1)$, $N(-1;-3)$. Phương trình tổng quát của  đường thẳng $\Delta$ đi qua $M$ và vuông góc với đường thẳng $MN$ là
	 \choice
		 {$4x-3y-5=0$}
		 {$2x+y+10=0$}
		{\True $3x+4y-10=0$} 
		 {$3x-4y-2=0$}
	 \loigiai{
		Đường thẳng $\Delta$ đi qua $M$ và vuông góc với đường  thẳng $MN$ nên nhận $\overrightarrow{NM}=(3;4)$ làm véc-tơ pháp tuyến.\\
		Phương trình đường thẳng $\Delta\colon 3(x-2)+4(y-1)=0\Leftrightarrow  3x+4y-10=0$.
	 }
\end{ex}

% nguồn: MaHoaNH_TracNghiem_Lop10
\begin{ex}%[0H4B1-2]%[ma:40959eaa]
	Phương trình tham số  của đường thẳng đi qua 2 điểm $A(3;-1)$ và $B(1; 5)$ là
	 \choice
		{\True  $\heva{&x=3+t \\ &y=-1-3 t}$}
		{$\heva{&x=3-t \\  &y=-1-3 t}$}
		{$\heva{&x=3+t  \\ &y=-1+3 t}$}
		{$\heva{&x=1-t \\  &y=5-3 t}$}
	 \loigiai{
		Ta có $\overrightarrow{AB}=(-2;6)$. Véctơ chỉ phương của  đường thẳng $AB$ là $\vec{u}=\dfrac{-1}{2}\overrightarrow{AB}=(1;-3)$.\\
		Phương trình tham số của đường thằng  $d$ là
		 \[d\colon \heva{&x=3+t \\ &y=-1-3 t.}\]
	 }
\end{ex}

% nguồn: MaHoaNH_TracNghiem_Lop10
\begin{ex}%[0H4B1-2]%[ma:826f800f]
	Đường thẳng đi qua $A(-1;2)$ nhận $\vv{n}=(2;4)$ làm véctơ pháp tuyến có phương trình là 
	 \choice
		{$x-2y-4=0$} 
		 {$x+y+4=0$}
		{\True  $x-2y+5=0$}
		 {$-x+2y-4=0$}
	 \loigiai{
		Phương trình đường  thẳng đi qua $A(-1;2)$ nhận $\vv{n}=(2;-4)$ làm véctơ pháp tuyến là
		\begin{equation*}
			2(x+1) - 4(y-2)=0\Leftrightarrow 2x-4y  +10=0\Leftrightarrow x-2y+5=0.
		\end{equation*}
	 }
\end{ex}

% nguồn: MaHoaNH_TracNghiem_Lop10
\begin{ex}%[0H4Y1-2]%[ma:29da36ac]
	Trong mặt phẳng $ Oxy$, đường thẳng đi  qua điểm $ M(2;-3)$, nhận véc-tơ $\overrightarrow{n}=\left(2;1\right)$ làm véc-tơ pháp tuyến có phương trình tổng quát là
	\choice 
		 {$ x+y+1=0$}
		{\True  $ 2x+y-1=0$}
		{$  2x+y-5=0$}
		 {$ 2x-3y-1=0$}
	\loigiai{Đường thẳng  đi qua điểm $ M(2;-3)$, nhận véc-tơ $\overrightarrow{n}=\left(2;1\right)$ làm véc-tơ pháp tuyến có phương trình tổng quát là
		\allowdisplaybreaks  \vspace*{-0.7cm}
		\begin{eqnarray*}
			& &  2(x-2)+1(y+3)=0 \\
			& \Leftrightarrow &  2x+y-1=0.
		\end{eqnarray*}
	} 
\end{ex}

% nguồn: MaHoaNH_TracNghiem_Lop10
\begin{ex}%[0H4B1-2]%[ma:f574a62a]
	Cho hình hành $ABCD$, biết $A(-2;1)$ và phương trình đường thẳng $CD$ là $3x-4y-5=0$. Phương trình tham số của đường thẳng $AB$ là 
	\choice 
		{\True $\heva{&x=-2-4t\\&y=1-3t}$} 
		{$\heva{&x=-2+3t\\&y=-2-2t}$} 
		 {$\heva{&x=-2-3t\\&y=1-4t}$}
		 {$\heva{&x=-2-3t\\&y=1+4t}$}
	\loigiai{ 
		Vì $ABCD$ là hình bình hành nên $AB$  song song với $CD$.\\
		Do đó, véctơ pháp tuyến của đường thẳng $CD$ cũng  là véctơ pháp tuyến của phương trình đường thẳng $AB$.\\
		Vậy $n_{\vv{AB}} =  (3;-4)$.\\
		Suy ra, véctơ chỉ phương của phương trình đường thẳng  $AB$ là $u_{\vv{AB}} = (4;3)=(-4;-3)$.\\
		Khi đó, phương  trình tham số của đường thẳng $AB$ là $\heva{&x=-2-4t\\&y=1-3t}$.
	} 
\end{ex}


%%%%%%%%%%%%%%%% ĐÚNG/SAI %%%%%%%%%%%%%%%%

% nguồn: MaHoaNH_DungSai_Lop10
\begin{ex}%[0H4N1-2]%[ma:15cf876f]
	Với mọi góc $\alpha$ thoả mãn $0^{\circ} \leq \alpha \leq 180^{\circ}$,  ta có
	 \choiceTF
		{\True $\sin \left(180^{\circ}-\alpha\right)=\sin \alpha$} 
		{$\cos \left(180^{\circ}-\alpha\right)=\cos  \alpha$}
		{$\tan \left(180^{\circ}-\alpha\right)=\tan  \alpha\left(\alpha \neq 90^{\circ}\right)$}
		{\True  $\cot \left(180^{\circ}-\alpha\right)=-\cot \alpha\left(0^{\circ}<\alpha<180^{\circ}\right)$}
	 \loigiai{
		Với mọi góc $\alpha$  thoả mãn $0^{\circ} \leq \alpha \leq 180^{\circ}$, ta có
		\begin{itemchoice}
			\itemch $\sin  \left(180^{\circ}-\alpha\right)=\sin \alpha$.
			 \itemch $\cos \left(180^{\circ}-\alpha\right)=-\cos \alpha$.
			\itemch $\tan \left(180^{\circ}-\alpha\right)=-\tan \alpha\left(\alpha \neq 90^{\circ}\right)$. 
			\itemch $\cot  \left(180^{\circ}-\alpha\right)=-\cot \alpha\left(0^{\circ}<\alpha<180^{\circ}\right)$.
		\end{itemchoice}
	 }
\end{ex}

% nguồn: MaHoaNH_DungSai_Lop10
\begin{ex}%[0H4H1-2]%[ma:be41416b]
	Cho góc $\alpha$ thỏa mãn  $\sin \alpha=\dfrac{4}{5}$, $ \left(90^\circ<\alpha<180^\circ\right)$. Khi đó
	\choiceTF 
		{\True $\cos \alpha=-\dfrac{3}{5}$} 
		{\True $\tan \alpha=-\dfrac{4}{3}$} 
		 {\True $\cot \alpha=-\dfrac{3}{4}$}
		{Giá trị  biểu thức $A=5 \sin ^{2} x+2 \cos ^{2} x$ bằng $\dfrac{62}{25}$}
	\loigiai{ 
		\begin{itemchoice}
			\itemch Ta có $\cos^2 \alpha=1-\sin^2\alpha=1-\left(\dfrac{4}{5}\right)^2=\dfrac{9}{25}$. Suy ra $\cos  \alpha =\pm \dfrac{3}{5}$. \\
			Vì $90^\circ<\alpha<180^\circ$  nên $\cos \alpha<0$. Vậy $\cos \alpha =-\dfrac{3}{5}$
			\itemch  Ta có $\tan \alpha=\dfrac{\sin \alpha}{\cos \alpha}=-\dfrac{4}{3}$.
			 \itemch Ta có $\cot \alpha=\dfrac{\cos \alpha}{\sin \alpha}=-\dfrac{3}{4}$.
			\itemch Ta  có $\sin \alpha=\dfrac{4}{5}\Rightarrow \sin^2x=\dfrac{16}{25}\Leftrightarrow \cos^2 x=\dfrac{9}{25}$.\\
			Thay vào biểu thức $A$ ta được $A=5\cdot\dfrac{16}{25}+2\cdot  \dfrac{9}{25}=\dfrac{98}{25}$.
		\end{itemchoice}
	} 
\end{ex}


%%%%%%%%%%%%%%%% TRẢ LỜI NGẮN %%%%%%%%%%%%%%%%

% nguồn: MaHoaNH_TraLoiNgan_Lop10
\begin{bt}%[0H4N1-2]%[ma:8f40fe28]
	Cho  $\alpha$ là góc nhọn và $\sin\alpha=\dfrac{5}{13}$. Tính $\cos\alpha$. (Làm tròn kết quả đến hàng phần trăm)
	\shortans{0,92} 
	\loigiai{ 
		 Ta có $\cos^{2}\alpha=1-\sin^{2}\alpha=1-\left(\dfrac{5}{13}\right)^{2}=\dfrac{144}{169}$.\\
		Do $\alpha$ là góc nhọn  nên $\cos\alpha>0$.\\
		 $\cos\alpha=\sqrt{\dfrac{144}{169}}=\dfrac{12}{13}\approx0{,}92$.
	} 
\end{bt}

% nguồn: MaHoaNH_TraLoiNgan_Lop10
\begin{bt}%[0H4V1-2]%[ma:fb13e31d]
	Cho $\cot\alpha=\dfrac{1}{2}$. Tính giá trị của  biểu thức $A=\dfrac{3\sin\alpha+4\cos\alpha}{2\sin\alpha-5\cos\alpha}$.
	\shortans{-10} 
	 \loigiai{
		Do  $\cot\alpha=\dfrac{\cos\alpha}{\sin\alpha}=\dfrac{1}{2}\Rightarrow\sin\alpha\ne 0$.\\
	Chia hai vế  biểu thức $A$ cho $\sin\alpha$, ta có $A=\dfrac{3+4\dfrac{\cos\alpha}{\sin\alpha}}{2-5\dfrac{\cos\alpha}{\sin\alpha}}=\dfrac{3+4\cot\alpha}{2-5\cot\alpha}=\dfrac{3+4\cdot\dfrac{1}{2}}{2-5\cdot\dfrac{1}{2}}=-10$.}
\end{bt}


%%%%%%%%%%%%%%%% TỰ LUẬN %%%%%%%%%%%%%%%%

% nguồn: MaHoaNH_TuLuan_Lop10
\begin{bt}%[0H4B1-2]%[0H4B1-4]%[ma:ef7ed958]
	[0H4B2-2] 
	Trong mặt phẳng tọa độ $Oxy$, cho hai điểm $A(-2;1)$, $B(2;4)$ và đường thẳng $d\colon3x+4y-10=0$. 
	\begin{enumerate}
		\item Viết phương trình đường thẳng $\Delta$ đi qua hai điểm $A(-2;1)$ và  $B(2;4)$.
		\item Tính cosin góc giữa hai đường thẳng  $d$ và $\Delta$.
		\item Viết phương trình  đường tròn đường kính $AB$.
	\end{enumerate}
	 \loigiai{
		\begin{enumerate}
			\item Đường thẳng  $AB$ nhận $\vv{AB}=(4;3)$ làm véctơ chỉ phương, do đó một véctơ pháp tuyến của đuởng thẳng $AB$ là $\vv{n}=(3;-4)$.\\
			Vậy phương trình tổng quát  của đường thẳng $AB$ là
			\begin{equation*}
				3(x+2)-4(y-1) = 0\Leftrightarrow  3x-4y+10=0.
			\end{equation*}
			\item Véctơ pháp tuyến của đường thẳng $d$ là $\vv{n}_d=(3;4)$,  $\vv{n}_{\Delta}=(3;-4)$.
			\begin{equation*}
				\cos(d,\Delta) = \left|\cos\left(\vv{n_d},\vv{n}_{\Delta}\right)\right|=\dfrac{\left|\vv{n}_d\cdot\vv{n}_\Delta\right|}{\left|\vv{n}_d\right|\cdot\left|\vv{n}_\Delta\right|}=\dfrac{7}{25}. 
			\end{equation*}
			\item Tâm $I\left(0;\dfrac{5}{2}\right)$, $R=\dfrac{5}{2}$.\\ 
			Phương trình đường tròn đường kính $AB$ là 
			\begin{equation*}
				x^2 + \left(y-\dfrac{5}{2}\right)^2=\dfrac{25}{4}. 
			\end{equation*}
		\end{enumerate}
	 }
\end{bt}

% nguồn: MaHoaNH_TuLuan_Lop10
\begin{bt}%[0H4K1-2]%[ma:bfedbf6b]
	Trong mặt phẳng $Oxy$, cho đường tròn $(C)\colon (x-3)^{2}+(y-4)^{2}=36$ và điểm $P(-3;-2)$. Từ điểm $P$ kẻ các tiếp tuyến $PM$ và $PN$  tới đường tròn $(C)$, với $M$, $N$ là các tiếp điểm. Viết phương trình tổng quát của đường thẳng $MN$.
	 \loigiai{
	 \immini{
		Đường tròn $(C)$ có tâm $I(3;4)$  và có bán kính $R=6$.\\
		Ta  có $MN\perp IP$ nên đường thẳng $MN$ có véc-tơ pháp tuyến là $\vec{IP}=(-6;-6)$ hay $MN$ có VTPT $\vec{n}=(1;1)$.\\
		Ta có $IM=R=6$, $PI=6\sqrt{2}$, ta suy ra tứ giác $IMPN$ là hình vuông. Gọi $K$ là trung điểm $IP$ $\Rightarrow  K(0;1)$. Ta có $MN$ qua $K(0;1)$.\\
		Phương trình tổng quát của đường  thẳng $MN$ là
		 \[
		 1\cdot (x-0)+1\cdot(y-1)=0\Leftrightarrow x+y-1=0.
		\] 
		}{\begin{tikzpicture}[scale=1,line join=round,line cap=round]
		\def\a{2.5} 
		\coordinate (I) at (0,0); 
		\draw (I) circle (\a); 
		\coordinate (M)  at ($(I)+(45:\a)$);
		\coordinate  (N) at ($(I)+(-45:\a)$);
		\path (I)--(M)--([turn]-90:\a) coordinate (Mt); 
		\draw (I)--(M)--(Mt); 
		\path (I)--(N)--([turn]90:\a)  coordinate (Nb);
		 \draw (I)--(N)--(Nb);
		\coordinate (P) at (intersection of M--Mt and N--Nb); 
		\draw  (P)--(M)
			(P)--(N) 
			 (M)--(N);
		\coordinate (K) at (intersection  of I--P and M--N);
		\draw (I)--(P); 
		\foreach  \x in {P,M,N}{
			\draw[fill=red] (\x) circle  (1pt) node[right]{$\x$};
		} 
		\foreach \x in {I}{ 
			\draw[fill=red] (\x) circle (1pt) node[left]{$\x$}; 
		 }
		\foreach \x  in {K}{
			\draw[fill=red] (\x)  circle (1pt) node[below right]{$\x$};
		 }
		\draw pic[draw=blue,angle  radius=2mm]
			 {right angle=I--M--P}
			pic[draw=blue,angle radius=2mm]  {right angle=I--N--P}
			 pic[draw=blue,angle radius=2mm] {right angle=M--K--P};
	\end{tikzpicture}
		} 
	} 
\end{bt}
